\begin{definition}[Conjugacy-class amplitudes]
    \label{def:peps_regular_flux_creation_vectors}
    \lean{TNLean.PEPS.regularFluxClassVector,TNLean.PEPS.normalizedRegularFluxClassVector,TNLean.PEPS.regularFluxPairVector}
    \leanok
    Let $G$ be a finite group and $C_G(g)$ the centralizer of $g$.
    Write $\kappa=|G|$ and $c_g=|C_G(g)|$. Define
    \begin{align}
        \alpha_g(t)  & =\sum_{z\in G}\mathbf 1_{t=zgz^{-1}},
                     & \widehat\alpha_g                      & =(\kappa c_g)^{-1/2}\alpha_g, \\
        \beta_g(a,b) & =\alpha_g(ab^{-1}).
    \end{align}
    These amplitudes retain the multiplicities in the accessible-register sum
    of \cite[Theorem 6.17]{Schuch2010PEPS}.
\end{definition}

\begin{theorem}[Exact conjugacy-sum normalization]
    \label{thm:peps_regular_flux_creation_norm}
    \lean{TNLean.PEPS.regularFluxClassVector_dotProduct,TNLean.PEPS.regularFluxClassVector_conjugation,TNLean.PEPS.regularFluxClassVector_one_eq_zero,TNLean.PEPS.normalizedRegularFluxClassVector_dotProduct,TNLean.PEPS.normalizedRegularFluxClassVector_orthogonal,TNLean.PEPS.normalizedRegularFluxClassVector_conjugation,TNLean.PEPS.normalizedRegularFluxClassVector_one,TNLean.PEPS.regularFluxPairVector_eq_sum,TNLean.PEPS.regularFluxPairVector_dotProduct}
    \leanok
    \uses{def:peps_regular_flux_creation_vectors}
    The squared norms are $\|\alpha_g\|^2=\kappa c_g$ and
    $\|\beta_g\|^2=\kappa^2c_g$. The normalized vector $\widehat\alpha_g$
    has norm one and is invariant under conjugation. For $g\ne1$ it is
    orthogonal to $\delta_1$, while $\widehat\alpha_1=\delta_1$.
    Moreover $\beta_g(a,b)=\sum_z\mathbf 1_{a=zgz^{-1}b}$.
\end{theorem}
\begin{proof}\leanok
    For every conjugate $zgz^{-1}$, precisely $c_g$ conjugating elements
    produce that conjugate. Expanding the squared norm therefore gives
    $\kappa c_g$. The reference coordinate $b$ contributes another factor
    $\kappa$ to the pair norm. Left multiplication of the conjugating variable
    proves conjugation invariance. A nontrivial class excludes the identity;
    the trivial class contains the identity with multiplicity $\kappa$.
\end{proof}

\begin{definition}[Creation on accessible registers]
    \label{def:peps_regular_flux_creation_matrix}
    \lean{TNLean.PEPS.regularFluxCreationMatrix,TNLean.PEPS.regularFluxPairCreationMatrix}
    \leanok
    \uses{def:peps_regular_flux_creation_vectors}
    Set $U_g=I-(\delta_1-\widehat\alpha_g)(\delta_1-\widehat\alpha_g)^\dagger$.
    On the pair $(a,b)$, let $V_g$ act as $U_g\otimes I$ in the coordinates
    $(ab^{-1},b)$.
\end{definition}

\begin{theorem}[Equivariant accessible-register creation]
    \label{thm:peps_regular_flux_creation_matrix}
    \lean{TNLean.PEPS.regularFluxCreationMatrix_mem_unitaryGroup,TNLean.PEPS.regularFluxCreationMatrix_mulVec_single,TNLean.PEPS.regularFluxCreationMatrix_commute_conjugation,TNLean.PEPS.regularFluxPairCreationMatrix_mem_unitaryGroup,TNLean.PEPS.regularFluxPairCreationMatrix_mulVec_diagonal}
    \leanok
    \uses{def:peps_regular_flux_creation_matrix,thm:peps_regular_flux_creation_norm,lem:peps_unitary_vector_swap}
    Both $U_g$ and $V_g$ are unitary. The operator $U_g$ commutes with every
    conjugation permutation and sends $\delta_1$ to $\widehat\alpha_g$.
    Consequently
    \begin{align}
        V_g\Bigl(\sum_r|r,r\rangle\Bigr)
        =(\kappa c_g)^{-1/2}\sum_{s,z}|zgz^{-1}s,s\rangle.
    \end{align}
    This is the normalized accessible-register operation in
    \cite[Theorem 6.17]{Schuch2010PEPS}; its action on a contracted PEPS
    region is a separate assertion.
\end{theorem}
\begin{proof}\leanok
    For $g\ne1$, reflect in the difference of the two orthogonal unit vectors.
    Every conjugation fixes both vectors and hence commutes with the reflection.
    For $g=1$ the reflection is the identity. Tensoring with the identity and
    reindexing preserves unitarity; summing the untouched reference coordinate
    reduces the pair action to $U_g\delta_1=\widehat\alpha_g$.
\end{proof}
